\pdfoutput=1
\documentclass[11pt]{amsart}

\usepackage{amsmath,amssymb,amsthm}
\usepackage[margin=1.25in]{geometry}
\usepackage{booktabs}
\usepackage[hidelinks]{hyperref}

\theoremstyle{plain}
\newtheorem{theorem}{Theorem}
\newtheorem{corollary}[theorem]{Corollary}
\newtheorem*{lemA}{Lemma A}
\newtheorem*{lemB}{Lemma B}
\newtheorem*{lemC}{Lemma C}
\theoremstyle{remark}
\newtheorem{remark}{Remark}

\newcommand{\Z}{\mathbb{Z}}
\newcommand{\F}{\mathbb{F}}
\DeclareMathOperator{\IM}{IM}

\title{An improved lower bound for the Furstenberg--S\'ark\"ozy problem}

\author{JD Jones}
\email{jjones@all-saints.org}

\subjclass[2020]{11B30 (Primary); 11B75, 05C70 (Secondary)}
\keywords{Furstenberg--S\'ark\"ozy theorem, square-difference-free sets,
Paley graph, induced matchings}

\begin{document}

\begin{abstract}
Let $D(N)$ be the largest size of a subset of $\{1,\dots,N\}$ no two of whose
elements differ by a nonzero perfect square. We prove
$\liminf_{N\to\infty} \log D(N)/\log N \ge \alpha_\infty
 = 0.753741541837329405\dots$,
where $\alpha_\infty$ is the value of an explicit closed-form optimization over
an eleven-block pool: nine of Krachun's Paley chains together with two
square-DAG certificates on the square-free composite moduli $235$ and $299$.
The previous record is Krachun's $\alpha_\star = 0.752796455874514\dots$, the
first bound past $3/4$. Only one lemma is new: a lift of square-DAGs on
square-free composite moduli; Krachun's gluing lemmas are used as published.
As a corollary, for prime $q$ the point--line incidence graph of $\F_q^2$
contains an induced matching of size $\gg_\varepsilon q^{1/2+\alpha_\infty-\varepsilon}$.
The full argument, including the numerical value, is formalized and
machine-checked in Lean~4.
\end{abstract}

\maketitle

\section{Introduction}

By the Furstenberg--S\'ark\"ozy theorem \cite{Furstenberg,Sarkozy}, a set of
integers of positive upper density contains two elements whose difference is a
perfect square. Quantitatively, let $D(N)$ denote the largest size of a subset
of $\{1,\dots,N\}$ no two of whose elements differ by a nonzero perfect
square; such sets are called \emph{square-difference-free} (SDF). The best
upper bound is due to Green and Sawhney \cite{GreenSawhney}:
$D(N) \ll N e^{-c\sqrt{\log N}}$. In the other direction, the history of the
lower-bound exponent is short:

\begin{center}
\begin{tabular}{lll}
\toprule
Year & Source & Exponent \\
\midrule
1984 & Ruzsa \cite{Ruzsa} & $0.733077$ \\
2008, 2015 & Beigel--Gasarch \cite{BeigelGasarch}; Lewko \cite{Lewko} (independently) & $0.733412$ \\
2026 & Krachun \cite{Krachun} & $0.752796\dots$ \\
2026 & this note & $0.753741\dots$ \\
\bottomrule
\end{tabular}
\end{center}

\begin{theorem}\label{thm:main}
$\displaystyle \liminf_{N\to\infty} \frac{\log D(N)}{\log N} \;\ge\;
\alpha_\infty \;=\; 0.753741541837329405\dots$
\emph{(nearest IEEE double $0.7537415418373294$)}, where $\alpha_\infty$ is
given by the explicit closed form \eqref{eq:alphainf} evaluated on the
eleven-block pool of Section~\ref{sec:pool}: nine Paley chains and two
square-DAG certificates on the composite moduli $235$ and $299$, lifted by
Lemma~A below, glued by Krachun's Lemma~5, and converted into integer sets by
his Lemma~4.
\end{theorem}

Equivalently, in pointwise form: for every $\varepsilon > 0$ there is an
$N_0(\varepsilon)$ with $D(N) \ge N^{\alpha_\infty - \varepsilon}$ for all
$N \ge N_0(\varepsilon)$. The statement is a liminf: the argument supplies no
constant $c$ with $D(N) \ge c N^{\alpha_\infty}$ for all $N$, and no such
claim is made here.

\begin{corollary}\label{cor:im}
For prime $q$, let $\IM(2,q)$ be the largest induced matching in the
point--line incidence graph of $\F_q^2$. Then for every $\varepsilon > 0$,
\[
\IM(2,q) \;\gg_\varepsilon\; q^{1/2 + \alpha_\infty - \varepsilon}.
\]
\end{corollary}

The corollary follows from Hunter--Pohoata--Verstra\"ete--Zhang
\cite[Prop.~2.3]{HPVZ}: a square-difference-free
$A \subseteq \{1,\dots,\lfloor q/10\rfloor\}$ yields an induced point--line
matching in $\F_q^2$ of size $\Omega(\sqrt{q}\,|A|)$; the implication is a
black box in $A$. The statement is for prime $q$ only; it is not made for
prime powers.

\subsection*{Context}
The classical constructions are the height-$1$ case --- a support with no
square difference at all, no ranks, used as alternating digits --- with
exponent $\tfrac12(1 + \log t/\log m)$. The record for that quantity is
$\tfrac12(1 + \log 12/\log 205) = 0.733412$, from $m = 205$, $t = 12$. For
prime $m$ the height-$1$ exponent never exceeds $3/4$: a
square-difference-free support mod $p$ is an independent set in the Paley
graph, of size at most $\sqrt{p}$ when $p \equiv 1 \pmod 4$, and is a single
point when $p \equiv 3 \pmod 4$ or $p = 2$. Krachun \cite{Krachun} crossed
$3/4$ by making height, rather than width, the quantity that pays for itself
in the CRT glue. This note keeps his glue and widens the supply of blocks:
square-DAGs on square-free composite moduli, where height can be strictly
smaller than width ($11 < 17$ and $12 < 19$ for the two certificates of
Section~\ref{sec:pool}).

\section{Definitions}\label{sec:defs}

\subsection*{The arc set}
For $m \ge 2$,
\[
Q_m := \{\, z^2 \bmod m : z \in \Z \,\} \setminus \{0\},
\]
the full image of squaring modulo $m$ with $0$ removed. For square-free $m$,
by the Chinese remainder theorem, $d \in Q_m \cup \{0\}$ if and only if
$d \bmod q$ is a square (possibly $0$) in $\F_q$ for every prime $q \mid m$.
Elements of $Q_m$ may be \emph{non-units} --- zero in some CRT coordinates
but not all. This is part of the definition, not an implementation choice:
Remark~\ref{rem:units} exhibits a counterexample showing that the unit-only
variant makes Lemma~A false. Sizes: $|Q_{235}| = 71$, $|Q_{299}| = 83$.

\subsection*{Square-DAGs and rankings}
A \emph{square-DAG} is a set $S \subseteq \Z/m\Z$ such that the digraph on
$S$ with arcs $x \to y$ whenever $y - x \in Q_m$ is acyclic. Acyclicity
forces antisymmetry, since a two-way pair is a $2$-cycle; no hypothesis on
$-1$ being a nonresidue is needed.

A \emph{ranking} of $S$ is a map $h_0 : S \to \{0,\dots,H_0-1\}$ with
$h_0(x) > h_0(y)$ for every arc $x \to y$; call $H_0$ the height of the
ranking. The reversed longest-path rank is such a map, and it realizes the
minimum $H_0$ over all rankings; that minimum value is the \emph{height} of
the DAG (the number of layers, equivalently the longest directed path counted
in vertices).

\subsection*{Ranked blocks}
A \emph{ranked block} is a triple $(C, h, H)$ with $C \subseteq \Z/P\Z$ and
$h : C \to \{0,\dots,H-1\}$ strictly decreasing along every nonzero square
difference mod $P$: if $x \ne y$ in $C$ and $y - x$ is a square mod $P$, then
$h(x) > h(y)$.

\subsection*{Paley chains}
A \emph{Paley chain} \cite[Def.~2]{Krachun} for a prime
$p \equiv 3 \pmod 4$ is a tuple $(s_0,\dots,s_{t-1})$ in $\F_p$ with
$s_b - s_a$ a nonzero quadratic residue for all $a < b$. A Paley chain is a
square-DAG whose digraph is a transitive tournament, so its height equals its
length $t$.

\subsection*{Valuations}
For a nonzero integer $d$, $v_m(d) := \max\{ j : m^j \mid d \}$; for
square-free $m = q_1 \cdots q_k$ this equals $\min_l v_{q_l}(d)$. For a
nonzero class $d$ in $\Z/m^{2e}\Z$, $v_m(d)$ is well defined and $< 2e$,
computed on any representative.

\section{The lifting lemma}\label{sec:lemA}

The following lemma is the only new ingredient of this note.

\begin{lemA}[composite even-digit lift]
Let $m \ge 2$ be square-free, let $S$ be a square-DAG on $\Z/m\Z$ with
ranking $h_0$ of height $H_0$ and $|S| = t$, and let $e \ge 1$. Put
\[
C(m,S,e) := \Bigl\{\, x = \sum_{j=0}^{2e-1} x_j m^j \;:\;
0 \le x_j < m,\ x_{2j} \in S \text{ for } 0 \le j \le e-1 \,\Bigr\}
\subseteq \Z/m^{2e}\Z,
\]
so that $|C(m,S,e)| = (mt)^e$, and
\[
h(x) := \sum_{j=0}^{e-1} h_0(x_{2j}) \, H_0^{\,e-1-j}
\;\in\; \{0,\dots,H_0^e - 1\}.
\]
Then for all $x \ne y$ in $C$ with $y - x \equiv z^2 \pmod{m^{2e}}$ for some
integer $z$ and $y - x \ne 0$, one has $h(x) > h(y)$. That is,
$(C, h, H_0^e)$ is a ranked block on $\Z/m^{2e}\Z$ of size $(mt)^e$.
\end{lemA}

\begin{proof}
\emph{Step 0 (setup).}
Let $x \ne y$ in $C$ and let $d := (y - x) \bmod m^{2e}$ satisfy
$d \equiv z^2 \pmod{m^{2e}}$ with $d \ne 0$. Let $r$ be the least base-$m$
position at which the digit strings of $x$ and $y$ differ. Both are strings
of $2e$ digits, so $r \le 2e-1$ automatically. As integers,
$y - x = \sum_{l \ge r} (y_l - x_l) m^l$ with $y_r \ne x_r$, so
$m^r \mid (y - x)$ and $m^{r+1} \nmid (y - x)$; the same holds for
$d = (y - x) + \kappa m^{2e}$ with $\kappa \in \{0,1\}$, since
$r + 1 \le 2e$ gives $m^{r+1} \mid m^{2e}$. Hence $m^r \mid d$,
$m^{r+1} \nmid d$, that is $r = v_m(d)$. Step~2 sharpens this to $r$ even
and $r \le 2e-2$.

\emph{Step 1 (leading digit, no borrows).}
Every term with $l > r$ contributes a multiple of $m$ to $(y-x)/m^r$, so
$(y-x)/m^r \equiv y_r - x_r \pmod m$; and
$d/m^r \equiv (y-x)/m^r \pmod{m^{2e-r}}$ because
$\kappa m^{2e}/m^r \equiv 0$. The leading digit is therefore
\[
\delta := (d/m^r) \bmod m \equiv y_r - x_r \pmod m, \qquad \delta \ne 0.
\]

\emph{Step 2 (the valuation of a nonzero square residue is even, and
$r \le 2e-2$).}
Fix a prime $q_l \mid m$ and set $v_l := v_{q_l}(z)$. Write
$d = z^2 + \lambda m^{2e}$ for an integer $\lambda$. Then:
\begin{itemize}
\item if $2v_l < 2e$, then
$v_{q_l}(z^2) = 2v_l < 2e \le v_{q_l}(\lambda m^{2e})$, so
$v_{q_l}(d) = 2v_l$, which is \emph{even}; call such a coordinate
\emph{uncapped};
\item if $2v_l \ge 2e$, then $v_{q_l}(d) \ge 2e$ and no parity information
is available; call such a coordinate \emph{capped}.
\end{itemize}
Since $d \ne 0$ in $\Z/m^{2e}\Z$, not every coordinate can have
$v_{q_l}(d) \ge 2e$, so at least one coordinate is uncapped, and
\[
r = v_m(d) = \min_l v_{q_l}(d) = \min_{\text{uncapped } l} 2v_l,
\]
a minimum of even numbers: the capped coordinates all have
$v_{q_l}(d) \ge 2e$ and cannot achieve the minimum, which is $< 2e$. Hence
$r = 2j$ is even with $0 \le j \le e-1$.

Square-freeness is essential here. For $m = \prod q_l^{a_l}$ with a repeated
prime factor, $v_m(d) = \min_l \lfloor v_{q_l}(d)/a_l \rfloor$ and evenness
fails: take $m = 9$ and $d = 9 = 3^2$, where $v_9(d) = 1$ is odd.

\emph{Step 3 (global divisibility $m^j \mid z$).}
Every coordinate satisfies $v_{q_l}(z^2) \ge 2j$. For uncapped $l$ this
reads $2v_l \ge r = 2j$, since otherwise $v_{q_l}(d) = 2v_l < r$ would
contradict $r$ being the minimum; for capped $l$ it holds because
$2v_l \ge 2e \ge 2j + 2$. Hence $v_{q_l}(z) \ge j$ for every $l$, and by
square-freeness $m^j \mid z$. Put $u := z/m^j \in \Z$.

\emph{Step 4 (the leading digit is a nonzero square mod $m$).}
From $d = z^2 + \lambda m^{2e} = m^{2j}(u^2 + \lambda m^{2e-2j})$ we get
\[
d/m^{2j} \equiv u^2 \pmod{m^{2e-2j}}, \qquad \text{with } 2e - 2j \ge 2.
\]
Reducing modulo $m$ and combining with Step~1: $\delta \equiv u^2 \pmod m$
and $\delta \ne 0$, so $\delta \in Q_m$. Here $u$ need not be a unit modulo
$m$, so $\delta$ may be a non-unit square, zero in some CRT coordinates.
This case occurs for genuine differences of the certificates of
Section~\ref{sec:pool}, which is why $Q_m$ must be the full image of
squaring.

\emph{Step 5 (rank drop).}
The first differing position $r = 2j$ is even, so $x_{2j}, y_{2j} \in S$,
and $\delta \equiv y_{2j} - x_{2j} \pmod m$ with $\delta \in Q_m$ gives an
arc $x_{2j} \to y_{2j}$, hence $h_0(x_{2j}) > h_0(y_{2j})$, a drop of at
least $1$. All even positions $2j' < 2j$ carry equal digits and contribute
equally to $h$. Therefore
\begin{align*}
h(x) - h(y)
&= \bigl[h_0(x_{2j}) - h_0(y_{2j})\bigr] H_0^{\,e-1-j}
 + \sum_{j' > j} \bigl[h_0(x_{2j'}) - h_0(y_{2j'})\bigr] H_0^{\,e-1-j'} \\
&\ge H_0^{\,e-1-j} - (H_0 - 1) \sum_{j' > j} H_0^{\,e-1-j'}
 = H_0^{\,e-1-j} - \bigl(H_0^{\,e-1-j} - 1\bigr) = 1 > 0. \qedhere
\end{align*}
\end{proof}

\begin{remark}[the unit-only variant is false]\label{rem:units}
Suppose the arc set were taken to be the unit squares mod $m$ only. Let
$m = 15$, whose unit squares are $\{1, 4\}$. Then $S = \{0, 1, 6\}$ is
acyclic for unit arcs --- the only such arc is $0 \to 1$ --- and
$h_0(0) = 1$, $h_0(1) = 0$, $h_0(6) = 1$ is a valid ranking for that arc
set. Lift with $e = 1$: $x = 0$ and $y = 36 = 6 + 2 \cdot 15$ both lie in
$C$, and $y - x = 36 = 6^2$ is a square modulo $225$ with $v_{15} = 0$ and
leading digit $6$, a non-unit square mod $15$ (the full square set mod $15$
is $\{1, 4, 6, 9, 10\}$). But $h(x) = 1 \le 1 = h(y)$, so the lifted
ranking fails, and Lemma~B applied to it would manufacture a false SDF set.
The definition $Q_m = $ image of squaring minus $0$ is therefore a
hypothesis of the lemma.
\end{remark}

\begin{remark}[what differs from the prime case]\label{rem:prime}
Four points, and nothing else.
(1) In Step~2, for prime $p$ the capped case collapses, since a capped
single coordinate forces $d \equiv 0$; for composite $m$ a coordinate can
hit the cap while the difference stays nonzero, so evenness must come from
the minimum over uncapped coordinates.
(2) In the prime case $p^j \mid z$ is immediate from
$v_p(d) = 2j = 2 v_p(z)$; for composite $m$ the divisibility must be
assembled coordinate by coordinate, the capped coordinates included. Step~3
is the composite form of the ``$n^{2j} \mid z^2$, hence $n^j \mid z$'' step
inside Krachun's Lemma~4.
(3) In Steps~4 and~5 the leading digit is automatically a unit square for
prime $p$, and may be a non-unit square for composite $m$; transplanting the
phrase ``nonzero quadratic residue modulo $p$'' verbatim would produce a
false lemma.
(4) Krachun extracts the order from his chain condition (2.1) --- every
forward difference along the chain a nonzero quadratic residue --- together
with $-1$ being a nonresidue mod $p$, which forces $b > a$; here that is
replaced by ``acyclicity yields a ranking,'' so $p \equiv 3 \pmod 4$ is not
needed in any form, and the height $t^e$ becomes $H_0^e$.
Everything else --- the digit representation, the borrow-free leading digit,
the lexicographic descent, the size count --- is Krachun's mechanism
verbatim.
\end{remark}

\section{Krachun's lemmas, used as published}\label{sec:krachun}

\begin{lemB}[{= \cite[Lemma~4]{Krachun}}]
Let $P = n^2$ be a perfect square, $n$ any positive integer, with no
primality or square-freeness hypothesis, and let $(C, h, H)$ be a ranked
block on $\Z/P\Z$. Then for every $L \ge 1$ there is a
square-difference-free $A_L \subseteq \{1,\dots,(PH)^L\}$ with
$|A_L| = |C|^L$.
\end{lemB}

The mechanism, in one sentence: base-$P$ words
$(x_0,\dots,x_{L-1}) \in C^L$ give $X = \sum \bar{x}_j P^j$, where
$\bar{x}_j \in \{0,\dots,P-1\}$ is the representative of $x_j$, with rank
$h_L(X) = \sum h(x_j) H^{L-1-j}$, and $B_L = \{X + P^L h_L(X)\}$ works,
because a positive square difference in $B_L$ would force $h_L$ not to
decrease along its own reduction mod $P^L$, contradicting the strict
decrease. For the proof see \cite[Lemma~4]{Krachun}. Our
$P = \prod m_i^{2e_i} = \bigl(\prod m_i^{e_i}\bigr)^2$ satisfies the
hypothesis; nothing needs extending.

\begin{lemC}[{= \cite[Lemma~5]{Krachun}}]
For $1 \le i \le \ell$ let $P_i$ be pairwise coprime perfect squares and
$(C_i, h_i, H_i)$ ranked blocks on $\Z/P_i\Z$. Under the CRT
identification, $C = \prod C_i \subseteq \Z/P\Z$ with $P = \prod P_i$ and
$h(x_1,\dots,x_\ell) = \sum h_i(x_i)$ is a ranked block of height
$H = 1 + \sum (H_i - 1)$.
\end{lemC}

The mechanism, in one sentence: the reduction of a square mod $P$ is a
square mod each $P_i$, possibly $0$; a zero coordinate leaves its rank
unchanged, a nonzero one drops it strictly, and at least one is nonzero.
For the proof see \cite[Lemma~5]{Krachun}. The statement never mentions
chains or primes, so it accepts arbitrary ranked blocks as given.

\section{The pool}\label{sec:pool}

Eleven blocks, all moduli pairwise coprime. Krachun's chain at $p = 23$ is
omitted because $23 \mid 299$.

\subsection{Nine Paley chains}
These are Krachun's chains \cite{Krachun}; for these,
$(m, t, H_0) = (p, t, t)$.

\begin{center}
\begin{tabular}{rlr}
\toprule
$p$ & chain & $t$ \\
\midrule
$3$   & $(0, 1)$ & $2$ \\
$7$   & $(0, 4, 1)$ & $3$ \\
$11$  & $(0, 3, 1, 4)$ & $4$ \\
$19$  & $(0, 5, 11, 9, 16)$ & $5$ \\
$31$  & $(0, 25, 14, 1, 19, 8, 2)$ & $7$ \\
$43$  & $(0, 31, 9, 23, 4, 40, 1)$ & $7$ \\
$59$  & $(0, 49, 15, 7, 16, 19, 35, 36, 5)$ & $9$ \\
$71$  & $(0, 8, 12, 18, 48, 27, 1, 37, 20)$ & $9$ \\
$103$ & $(0, 79, 25, 58, 55, 81, 4, 1, 34, 83, 59)$ & $11$ \\
\bottomrule
\end{tabular}
\end{center}

Krachun's tenth chain, $23 : (0, 18, 1, 3, 4)$ with $t = 5$, is used here
only to restate his constant $\alpha_\star$.

\subsection{Composite block 1:
\texorpdfstring{$(m, t, H_0) = (235, 17, 11)$, with $235 = 5 \cdot 47$
square-free}{(m, t, H0) = (235, 17, 11), with 235 = 5.47 square-free}}

Support with its exhibited ranking, written as \emph{vertex\,: rank}:

\begin{center}
\begin{tabular}{cccccc}
$0 : 10$ & $112 : 10$ & $196 : 9$ & $224 : 9$ & $155 : 8$ & $136 : 7$ \\
$67 : 6$ & $110 : 6$  & $92 : 5$  & $189 : 5$ & $126 : 4$ & $193 : 4$ \\
$22 : 3$ & $50 : 3$   & $64 : 2$  & $73 : 1$  & $148 : 0$ & \\
\end{tabular}
\end{center}

Seventeen vertices, ranks $0,\dots,10$, so height $11$. Checking this is a
square-DAG with this ranking is one pass over the $17 \cdot 16$ ordered
pairs: for each pair test whether the difference lies in $Q_{235}$, and if
so that the rank strictly drops.

\subsection{Composite block 2:
\texorpdfstring{$(m, t, H_0) = (299, 19, 12)$, with $299 = 13 \cdot 23$
square-free}{(m, t, H0) = (299, 19, 12), with 299 = 13.23 square-free}}

Support:
\[
7,\ 21,\ 25,\ 40,\ 46,\ 78,\ 83,\ 116,\ 153,\ 161,\ 165,\ 206,\ 207,\ 210,\
212,\ 244,\ 264,\ 289,\ 292.
\]
Nineteen vertices. The induced digraph under $Q_{299}$ is acyclic, and its
reversed longest-path rank has height $12$; both are computed directly from
the support in milliseconds.

Both blocks are presented as lower-bound certificates only. The cited
heights $H_0 = 11$ and $12$ are the heights of these exhibited certificates
and nothing more; no minimality, maximality, or census claim is made about
them or about the pool.

\section{The exponent}\label{sec:exponent}

Write $m_i$, $t_i$, $H_i$ for the modulus, size, and height $H_0$ of pool
block $i$; after Lemma~A with multiplicity $e_i$, block $i$ is a ranked
block of height $H_i^{e_i}$, which is what Lemma~C receives. Glue across
the pool with Lemma~C and feed the result to Lemma~B at length $L$. Then
\[
P = \prod_i m_i^{2e_i}, \qquad
|C| = \prod_i (m_i t_i)^{e_i}, \qquad
H = 1 + \sum_i \bigl(H_i^{e_i} - 1\bigr),
\]
and the SDF sets produced by Lemma~B have exponent
\begin{equation}\label{eq:alphae}
\alpha(e) = \frac{\sum_i e_i \log(m_i t_i)}
{2 \sum_i e_i \log m_i + \log\bigl(1 + \sum_i (H_i^{e_i} - 1)\bigr)}.
\end{equation}

\subsection*{The limit}
Every pool block has $H_i \ge 2$. Allocate
$e_i(U) := \lfloor U / \log H_i \rfloor$ for a real parameter
$U \ge \max_i \log H_i$, so that every $e_i(U) \ge 1$ and Lemma~A applies.
Then $e_i(U) \log H_i \in (U - \log H_i,\, U]$, so
$e_i(U) = U/\log H_i + O(1)$, and
\[
\max_i H_i^{e_i} \le 1 + \sum_i \bigl(H_i^{e_i} - 1\bigr)
\le \ell \cdot \max_i H_i^{e_i}
\quad\Longrightarrow\quad \log H = U + O(1),
\]
with implied constants depending only on the pool. Substituting,
\[
\text{numerator} = U \sum_i \frac{\log(m_i t_i)}{\log H_i} + O(1),
\qquad
\text{denominator} = U \Bigl(1 + 2 \sum_i \frac{\log m_i}{\log H_i}\Bigr)
+ O(1),
\]
both $\Theta(U)$, so $\alpha(e(U)) \to \alpha_\infty$ as $U \to \infty$,
where
\begin{equation}\label{eq:alphainf}
\alpha_\infty
= \frac{\displaystyle \sum_i \frac{\log(m_i t_i)}{\log H_i}}
       {\displaystyle 1 + 2 \sum_i \frac{\log m_i}{\log H_i}}.
\end{equation}
This is Krachun's Theorem~1 computation with $\log t_i$ replaced by
$\log H_i$; for a Paley chain $H = t$, so his formula is the special case.

\subsection*{Passage to all \texorpdfstring{$N$}{N}}
Fix $\rho < \alpha_\infty$ and choose $U$ with
$\alpha := \alpha(e(U)) > \rho$. Write $P = P(U)$, $H = H(U)$,
$N_L := (PH)^L$, so Lemma~B gives SDF $A_L \subseteq \{1,\dots,N_L\}$ with
$|A_L| = |C|^L = N_L^\alpha$. For arbitrary $N \ge PH$ set
$L := \lfloor \log N / \log(PH) \rfloor \ge 1$. Then $N_L \le N$ and
$A_L \subseteq \{1,\dots,N\}$ is square-difference-free, so
\[
D(N) \ge |A_L| = N_L^\alpha > \bigl(N/(PH)\bigr)^\alpha
= (PH)^{-\alpha} N^\alpha,
\]
whence
\[
\frac{\log D(N)}{\log N} \ge \alpha - \alpha \frac{\log(PH)}{\log N}
\longrightarrow \alpha > \rho,
\]
so $\liminf_{N\to\infty} \log D(N)/\log N \ge \rho$. Letting
$\rho \nearrow \alpha_\infty$ proves Theorem~\ref{thm:main}. \qed

\subsection*{Finite-stage certifiability}
No limit is needed for an unconditional bound. Each finite $U$ certifies
its own exponent with an explicit constant: with $P = P(U)$ and $H = H(U)$,
\[
D(N) \ge (PH)^{-\alpha(e(U))} \, N^{\alpha(e(U))}
\qquad \text{for all } N \ge PH.
\]
The limit is approached from below through these finite constructions:
$\alpha(e(U)) = 0.744217$ at $U = 5$, $0.751542$ at $U = 20$, $0.753210$
at $U = 80$, $0.753616$ at $U = 320$, gap of order $1/U$.

\subsection*{The value}
Evaluating \eqref{eq:alphainf} on the eleven-block pool:
\[
\alpha_\infty = 0.753741541837329405\dots
\qquad \text{(nearest double } 0.7537415418373294\text{)}.
\]
For comparison, the same closed form over Krachun's ten Paley pairs
(his (1.1)) gives $\alpha_\star = 0.752796455874514\dots$

\subsection*{Scope}
No optimality is claimed --- not of the two certificates, not of their
heights, not of the pool, and not of the method. The blocks are lower-bound
certificates: exhibited objects that a reader checks in milliseconds.
Krachun's \S3 discussion of the limit of the approach is heuristic,
restricted to prime moduli, and stated by its author to be without precise
proofs, so it bears on none of the above.

\section{Verification and formalization}\label{sec:verification}

The finite content of the note --- the arc sets $Q_{235}$ and $Q_{299}$
rebuilt from the definition, both certificates checked for size, acyclicity,
and longest-path height, all ten Paley chains checked for primality,
$p \equiv 3 \pmod 4$, distinctness, and every forward difference a nonzero
quadratic residue, and both constants recomputed in high-precision
decimal --- is re-checked by the ancillary script \texttt{verify.py}
(Python~3, standard library only, well under a minute, all data inlined,
nonzero exit status on any failure). The flag \texttt{--lift-demo}
additionally builds the $e = 1$, $L = 1$ lifted integer set for the $235$
block and brute-checks square-difference-freeness --- an empirical witness
of Lemmas~A and~B at the smallest scale, a demonstration, not a proof. The
script re-checks data stated in full above; it is not an input to the
mathematics.

The complete argument is moreover formalized in Lean~4 \cite{Lean4} with
Mathlib \cite{Mathlib}: the definitions of Section~\ref{sec:defs}, Lemma~A,
Lemmas~B and~C reproved rather than cited, the exponent computation, and
Theorem~\ref{thm:main} in both liminf and pointwise $\varepsilon$-form,
together with the numerical statement $0.7537 < \alpha_\infty$, certified by
kernel-checked integer comparisons with no floating-point or transcendental
arithmetic anywhere. The development comprises $126$ declarations, every one
depending only on Lean's three standard axioms (\texttt{propext},
\texttt{Classical.choice}, \texttt{Quot.sound}). A source snapshot of the
formalization accompanies this submission as ancillary files
(\texttt{anc/lean/}), with the toolchain and dependency versions pinned by
\texttt{lean-toolchain} and \texttt{lake-manifest.json}, so the kernel check
can be reproduced against the exact versions used. The formalization is
maintained at \url{https://github.com/JD-Jones-ASES/fs-lower-bound-lean};
the note's repository, including the verification script, is at
\url{https://github.com/JD-Jones-ASES/fs-lower-bound}.

\section*{Provenance}

This note reports AI-generated mathematics with a human managing the
workflow. The construction, proofs, and text were produced by Claude Code
(Fable~5, Anthropic), with contributions from Codex (GPT~5.6 Sol, OpenAI)
and Grok Build (Grok~4.6, xAI); the author directed the work and takes
responsibility for the content. Verification does not require trusting any
of these systems: everything needed to reimplement the construction is
stated above, the finite data are re-checked by the ancillary script, and
the complete argument is machine-checked in Lean~4
(Section~\ref{sec:verification}).

\begin{thebibliography}{9}

\bibitem{BeigelGasarch}
R.~Beigel and W.~Gasarch,
\emph{Square-difference-free sets of size $\Omega(n^{0.7334})$},
arXiv:0804.4892 (2008).

\bibitem{Furstenberg}
H.~Furstenberg,
\emph{Ergodic behavior of diagonal measures and a theorem of Szemer\'edi on
arithmetic progressions},
J. Analyse Math. \textbf{31} (1977), 204--256.

\bibitem{GreenSawhney}
B.~Green and M.~Sawhney,
\emph{New bounds for the Furstenberg--S\'ark\"ozy theorem},
arXiv:2411.17448 (2024).

\bibitem{HPVZ}
Z.~Hunter, C.~Pohoata, J.~Verstra\"ete, and S.~Zhang,
\emph{Large point-line matchings and small Nikodym sets},
arXiv:2601.19879 (2026).

\bibitem{Krachun}
D.~Krachun,
\emph{Square-difference-free sets beyond the three-quarter barrier},
arXiv:2608.01325 (2026).

\bibitem{Lean4}
L.~de~Moura and S.~Ullrich,
\emph{The Lean 4 theorem prover and programming language},
in: Automated Deduction -- CADE 28, Lecture Notes in Comput. Sci.
\textbf{12699}, Springer, 2021, 625--635.

\bibitem{Lewko}
M.~Lewko,
\emph{An improved lower bound related to the Furstenberg--S\'ark\"ozy
theorem},
Electron. J. Combin. \textbf{22} (2015), \#P1.32.

\bibitem{Mathlib}
The mathlib Community,
\emph{The Lean mathematical library},
in: Proc. 9th ACM SIGPLAN Int. Conf. on Certified Programs and Proofs
(CPP 2020), ACM, 2020, 367--381.

\bibitem{Ruzsa}
I.~Z.~Ruzsa,
\emph{Difference sets without squares},
Period. Math. Hungar. \textbf{15} (1984), 205--209.

\bibitem{Sarkozy}
A.~S\'ark\"ozy,
\emph{On difference sets of sequences of integers I},
Acta Math. Acad. Sci. Hungar. \textbf{31} (1978), 125--149.

\end{thebibliography}

\end{document}
